\chapter{磁场}
\section{安培定律}

人们对于磁学的研究最早来源于天然磁体对于铁磁性物质的吸引和磁体间的相互作用. 对于一个条形磁体, 其两极吸引铁磁性物质的能力最强, 由其在地磁场作用下产生的指向, 人们将磁极区分为北磁极和南磁极.

磁极的同性相斥, 异性相吸的特点与电现象十分相似. 但早期人们对于磁和电的研究彼此独立地进行着, 直到奥斯特通过实验揭示了电与磁之间的联系.

奥斯特实验可以概述如下: 对于一根南北放置的导线, 在其下方放置一可在水平面内自由转动的磁针. 当导线中没有电流通过时, 磁针沿着地磁场方向沿南北取向. 但当导线中通入电流时, 磁针会发生偏转, 且电流反向时, 磁针的偏转方向也相反.

这一实验表明, 电流可以对磁铁施加作用力. 我们将一段水平的直导线悬挂在马蹄形磁铁的两极间并通入电流后, 导线会发生移动. 这表明磁铁也会对电流施加作用力. 相关的实验表明, 磁铁和磁铁, 电流和电流之间也有存在着相互作用力.

在电学的研究中, 电荷之间的作用力是通过电场来传递的. 电场的基本性质是会对置于其中的其他电荷施加作用力. 而磁极和电流间的作用也是这样的, 只不过是通过磁场来传递. 磁场的基本性质是会对任何置于其中的磁极或者电流施加作用力.

实验表明, 一个载流线圈的行为非常像一块磁铁. 安培在这一实验现象的基础上, 提出了这样的假说. 组成磁铁的最小单元(磁分子)就是环形电流. 若这样一些分子环流定向地排列起来, 在宏观上就会显示出$N,S$极出来. 以上的观点被称为安培分子环流假说. 那个时候人们尚不清楚原子的结构, 因而无法解释物质内部的分子环流是怎样形成的. 现在我们知道原子是由带正电的原子核和绕核旋转的负电子组成的. 电子不仅绕核旋转, 而且还有自旋, 这些运动形成了所谓 ``分子环流'', 成为了物质磁性的来源.

按照上面的观点, 磁铁的本质实际上是电荷的运动, 上面提到的各种相互作用都可以归结电流之间的相互作用, 这种相互作用是通过磁场来传递的.

因此我们现在来研究电流与电流间的相互作用. 然而在研究这一规律之前, 由于恒定电流只能存在于闭合回路中, 想要获得孤立的恒定电流是不可能的. 我们的基本方法是把两个相互作用的载流回路切分为许多无穷小的线元, 也就是电流元. 只要知道了任意一对电流元的作用规律, 整个闭合回路受的力就可以通过矢量叠加计算出来. 由于实验中无法实现一个孤立的恒定电流元, 无法直接用实验测定相互作用, 只能间接地从闭合载流回路的实验中倒推出来. 安培天才地设计了四个实验, 得到了最终的结果.

这里我们不看具体的推导过程, 直接来看结果.

\begin{theorem}[安培定律]

对于两个电流元$I_1\d \mathbf{l}_1$和$I_2\d \mathbf{l}_2$, 电流元$1$给电流元$2$的作用力可以写为:
\[
\d \mathbf{F}_{12} = k\frac{I_1I_2\d \mathbf{l}_2\times(\d \mathbf{l}_1\times {\hat{\mathbf{r}}_{12}})}{r^2_{12}}
\]
\end{theorem}

对这个式子我们做些说明: 考虑图示的情况.
  \insertfig{2-1.png}{电流元相互作用一定满足牛顿第三定律吗?}{0.25}

注意到这种情况下$1$对$2$的作用力为0, 但$2$对$1$的作用力不为0. 也就是说, 按照前面式子确定的电流元之间的作用力是不一定满足牛顿第三定律的. 然而实际中并不存在孤立的恒定电流元, 并且我们可以证明, 两个闭合回路间的总作用力是满足牛顿第三定律的.

电流元之间的作用力不满足牛顿第三定律, 表明动量和并不仅仅为电流或者质点所具有. 电场和磁场也可以携带动量与角动量, 而且电荷的动量与角动量的变化, 恰好由电磁场的角动量变化予以补偿, 从而维持了动量守恒和角动量守恒.

国际上现行的电磁学单位制是MKSA制, 除了长度, 质量, 时间之外的第四个基本量是电流, 其单位被定义为安培. 安培的定义和绝对测量正是以安培定律为依据的.

在安培定律中, 我们把比例系数$k$写作$\frac{\mu_0}{4\pi}$, 并取$\mu_0=4\pi\times10^{-7}$. 这样确定下来的单位即为安培.


\section{\texorpdfstring{磁感应强度\, 毕奥-萨伐尔定律}{磁感应强度 毕奥-萨伐尔定律}}
为了定量地描述电场的分布, 我们曾经引入过电场强度矢量的概念. 同样地, 为了描述电流间的相互作用, 我们需要定义一个描述磁场的矢量---磁感应强度$B$.

在磁场情形中, 与库仑定律相当的基本定律是安培定律, 也就是:
\[
\d \mathbf{F}_{12} = \frac{\mu_0}{4\pi} \frac{I_1I_2\d \mathbf{l}_2\times(\d \mathbf{l}_1\times {\hat{\mathbf{r}}_{12}})}{r^2_{12}}
\]

仿照电场形式, 我们把上式拆成两个部分:
\[
\d \mathbf{F}_{12}=I_2\d \mathbf{l}_2\times \d \mathbf{B}
\]
\[
\d \mathbf{B} = \frac{\mu_0}{4\pi}\frac{I_1 \d \mathbf{l}_1\times \hat{\mathbf{r}}_{12}}{r^2_{12}}
\]

$\d \mathbf{l}_1$本来是某个闭合回路的一部分, 整个回路$L_1$对于试探电流元$I_2 \d \mathbf{l}_2$的作用力是上式对$\d l_1$的积分:
\[
\d \mathbf{F}_2 = \frac{\mu_0}{4\pi}\oint_{L_1}\frac{I_2\d \mathbf{l}_2\times(\d \mathbf{l}_1\times {\hat{\mathbf{r}}_{12}})}{r^2_{12}}=\frac{\mu_0}{4\pi}I_2\d \mathbf{l}_2\times\oint_{L_1}\frac{I_1\d \mathbf{l}_1\times {\hat{\mathbf{r}}_{12}}}{r^2_{12}}
\]

于是按照上式, 我们定义:
\[
\d \mathbf{F}_2 = I_2\d \mathbf{l}_2\times \mathbf{B}, \; \mathbf{B} = \frac{\mu_0}{4\pi}\oint_{L_1}\frac{I_1 \d \mathbf{l}_1\times {\hat{\mathbf{r}}_{12}}}{r^2_{12}}
\]

像这样定义的$B$被称为磁感应强度矢量.

把$1$代表某个闭合载流回路的一个任意电流元, $2$代表任意场点并且略去下标不写, 那么就有:
\begin{theorem}[毕奥-萨伐尔定律]
\[
\d \mathbf{B}=  \frac{\mu_0}{4\pi}
\frac{I \d \mathbf{l}\times {\hat{\mathbf{r}}}}{r^2}
\]
\[
\mathbf{B} = \frac{\mu_0}{4\pi}
\oint_L \frac{I \d \mathbf{l}\times {\hat{\mathbf{r}}}}{r^2}
\]
\end{theorem}

有了毕奥-萨伐尔定律, 我们可以计算一些电流体系的磁感应强度.

\begin{example}
  \insertfig{2-2.png}{载流直导线的磁场}{0.25}
考虑一段直导线旁任意一点$P$的磁感应强度. 根据毕奥-萨伐尔定律可以看出任意电流元在该点产生的元磁场都是在$P$点垂直纸面向内的. 因此我们只需要计算$\d B$的代数和. 列出积分式:
\[
B = \frac{\mu_0}{4\pi}\int_{A_1}^{A_2}\frac{I\sin\theta\d l}{r^2}
\]

几何关系告诉我们:
\[
l=-r\cos\theta,r_0=r\sin\theta\quad\Rightarrow\quad l=-r_0\cot\theta,\d l = \frac{r_0 \d\theta}{\sin^2\theta}
\]
完成积分:
\[
B = \frac{\mu_0 I}{4\pi}(\cos\theta_1-\cos\theta_2)
\]
其中$\theta_1,\theta_2$分别是$A_1,A_2$与场点连线与导线电流的夹角.

对无限长直导线, $\theta_1=0,\theta_2=\pi$, 于是:
\[
B = \frac{\mu_0}{4\pi}\int_{-\infty}^{+\infty}\frac{I\sin\theta\d l}{r^2}=\frac{\mu_0}{4\pi}\int_0^\pi\frac{I}{r_0}\sin \theta\d\theta=\frac{\mu_0 I}{2\pi r_0}
\]
以上结果表明, 在载流无限长导线周围的磁感应强度大小与距离的一次方成反比.

\end{example}

\begin{example}
设圆线圈的中心为$O$, 半径为$R$. 其上任意点$A$处的电流元在对称轴线上一点$P$产生的元磁场方向如图所示, 根据对称性我们只需要计算沿着轴线方向的分量. 于是总体的磁场强度为:

\insertfig{2-3.png}{圆线圈轴线上的磁场}{0.25}
\[
B=\oint\d B \cos \alpha =\oint \frac{\mu_0}{4\pi} \frac{I\d l}{r_0^2}\sin^2\alpha\cos\alpha = \frac{\mu_0}{4\pi}\frac{IR}{(R^2+r_0^2)^{\frac{3}{2}}}\oint \d l =\frac{\mu_0R^2I}{2(R^2+r_0^2)^{\frac{3}{2}}}
\]
于是圆心处:
\[
B=\frac{\mu_0 I}{2R}
\]
远场条件:($r_0 \gg R$)
\[
B=\frac{\mu_0R^2I}{2r_0^3}
\]
\end{example}

利用这一结果, 我们来讨论一中更常见的产生近似匀强磁场分布的装置---亥姆霍兹线圈.

\begin{example}
一对相同的圆形线圈彼此平行共轴放置. 设两线圈内的电流都是$I$, 且回绕方向一致, 半径为$R$, 两者的间距为$a$. 求:

(1)求轴线上的磁场分布;

(2)求$a$为何值两线圈轴线上中点$O$处的磁场最均匀.
\insertfig{2-4.png}{亥姆霍兹线圈}{0.25}

利用前面的结果, 令$r_0=x\pm\frac{a}{2}$. 其中$x$为场点在轴线上的坐标.
则:
\[
B_1=\frac{\mu_0R^2I}{2[R^2+(x+\frac{a}{2})^2]^{\frac{3}{2}}}, \, B_2=\frac{\mu_0R^2I}{2[R^2+(x-\frac{a}{2})^2]^{\frac{3}{2}}}
\]

$B_1$,$B_2$两者方向相同, 于是$B=B_1+B_2$. 对$B$求一阶, 二阶导数, 得到:
\[
\frac{\d B}{\d x}=-\frac{3\mu_0R^2I}{2}\left\{ \frac{x+\frac{a}{2}}{[R^2+(x+\frac{a}{2})^2]^{\frac{5}{2}}}+\frac{x-\frac{a}{2}}{[R^2+(x-\frac{a}{2})^2]^{\frac{5}{2}}}\right\}
\]
这在$x=0$时恒为$0$.
\[
\frac{\d^2 B}{\d x^2} = \frac{3\mu_0R^2I}{2}\left\{ \frac{4(x+\frac{a}{2})^2-R^2}{[R^2+(x+\frac{a}{2})^2]^{\frac{7}{2}}}+\frac{4(x-\frac{a}{2})^2-R^2}{[R^2+(x-\frac{a}{2})^2]^{\frac{7}{2}}}\right\}
\]

为使二阶导也为$0$, 则$a=R$. 由于$B(x)$是偶函数, 其三阶导数在$0$处也为$0$. 于是这一装置可以在相当大的范围内维持均匀. 这样的半径等于间距的一对共轴线圈就被称为亥姆霍兹线圈.
\end{example}

\begin{example}
  绕在圆柱面上的螺线形线圈被称为螺线管. 当绕上的导线很细且是一匝一匝密绕的时候, 我们可以把其看作一个导体圆筒, 电流连续地沿着环向分布. 设其半径为$R$, 总长度为$l$, 单位长度内的电流为$i$. 取圆筒的轴线为$x$轴.

  \insertfig{2-5.png}{圆筒轴线上的磁场}{0.3}

  则在图中的P点处, 对于每一匝线圈, 其元磁场可以按照这样的方式计算:
  \[
  \d B = \frac{\mu_0R^2\d l\, i}{2}\frac{1}{[R^2+(x-l)^2]^\frac{3}{2}}
  \]

  取$r=\frac{R}{\sin\beta}$,$x-l=r\cos\beta$于是我们得到:
  \[
  \frac{x-l}{R}=\cot\beta, \,\frac{\d l}{R}=\frac{\d \beta}{\sin^2\beta}
  \]
  做变量替换, 完成积分:
  \[
  B=\frac{\mu_0 i}{2}\int_{\beta_1}^{\beta_2}\sin\beta\d\beta=\frac{\mu_0i}{2}(\cos\beta_1-\cos\beta_2)
  \]

  考虑特殊情形, 当圆筒无限长时,$\beta_1=0,\beta_2=\pi$
  \[
  B=\mu_0 i
  \]
  这表明在无限长环向电流圆筒内轴线上磁场均匀.

  对于半无限长圆筒的一端, $\beta_1=0,\beta_2=\frac{\pi}{2}$或者$\beta_1=\frac{\pi}{2},\beta_2=\pi$, 均有
  \[
  B = \frac{\mu_0i}{2}
  \]

  这结果是自然的, 因为我们可以将一个无限长圆筒从任意地方结成两半, 这两半对于总磁感应强度的贡献应该是一致的, 于是得到每一半的贡献.

  对于螺线管, 设其单位长度内的匝数为$n$, 每匝的电流为$I$, 于是$i=nI$, 做这样的替换之后, 就可以适用上面的公式.

  对于无穷大载流平面产生的磁场, 请参考5.4节中的推导.
\end{example}

\section{安培环路定理}

我们不加证明地给出安培环路定理.
\begin{theorem}[安培环路定理]

磁感应强度沿任何闭合环路$L$的线积分, 等于穿过这环路所有电流代数和的$\mu_0$倍. 也就是:
\[
\oint \mathbf{B} \cdot \d l = \mu_0 \sum_{enc} I
\]

\end{theorem}

其中电流的正负如下规定: 当穿过回路的电流方向与回路的环绕方向满足右手法则时电流为正, 反之为负.

利用安培环路定理我们可以计算一些对称性载流导线的磁场分布.

\begin{example}
考虑一根圆截面的无限长载流直导线, 其导线半径为$R$, 电流$I$均匀地通过横截面. 求其磁场分布.

我们取一个环绕着导线的环路, 如图所示.
\insertfig{2-6.png}{无限长圆截面直导线的磁场分布}{0.2}

(1)$r>R$时, 全部电流都通过了环路, 有:
\[
2\pi r B = \mu_0 I\Rightarrow B = \frac{\mu_0 I}{2\pi r}
\]

(2)$r<R$时, 有:
\[
2\pi r B = \mu_0 I \frac{r^2}{R^2} \Rightarrow B = \frac{\mu_0 Ir}{2\pi R^2}
\]
\end{example}

\begin{example}
对无限长密绕螺线管, 我们可以重新给出一个定理.
\insertfig{2-7.png}{无限长密绕螺线管的磁场分布}{0.25}

令n为单位长度上的线圈匝数, 在无穷远处, 磁感应强度为$0$. 于是我们取出图示的环路. 在无穷远处的路径积分为$0$. 伸向无穷远处的两条路径根据对称性总贡献也为$0$. 于是可能产生贡献的路径只剩下经过$P$的路径.

当$P$在螺线管外时, 通过该环路的电流为0, 于是$B=0$; 当$P$在螺线管内时, 通过的电流为$nI\Delta l$. 于是:
\[
B \Delta l = \mu_0 n I \Delta l \Rightarrow B = \mu_0 nI
\]

类似的, 绕在环面的螺形线圈我们有类似的结论. 令n为螺绕环单位长度上的线圈匝数, R为螺绕环的中心半径 (从环的中心到管状部分中心的距离).

\insertfig{2-8.png}{螺绕环的磁场}{0.25}

我们在环面内沿着半径为$r$的路线积分, 有:\[
2\pi Br =\mu_0 n I \cdot 2\pi R \Rightarrow B = \mu_0 nI
\]

\end{example}

\section{磁场的高斯定理}
仿照电通量, 我们规定通过一个曲面$S$的磁感应通量, 简称磁通量为:
\[
\Phi_B= \oiint_S \mathbf{B}\cdot \d S
\]
载流导线产生的磁感应线是闭合的, 于是直觉告诉我们从一个闭合曲面$S$穿入的磁感应线必然从另一处穿出, 于是通过任意闭合曲面的磁通量恒等于$0$. 也就是:
\begin{theorem}[磁场高斯定理]


\[
\oiint_S \mathbf{B}\cdot \d S =0
\]
这结论被称为磁场的高斯定理.
\end{theorem}

我们从毕奥-萨伐尔定律出发证明这点.
\insertfig{2-9.png}{证明磁场高斯定理}{0.25}

考虑单个电流元, 其产生的磁感应线是以$\d l$方向为轴线的圆, 且在这个圆周上磁场强度大小处处相等.

取一个闭合曲面$S$. 每根圆形的闭合磁感线或者不与$S$相交, 或者穿过$S$偶数次. 我们取一个闭合的以$\d l$方向为轴线的环形管, 则穿过该环形管内的任意截面磁通量为常量, 因为我们都可以通过投影面积的办法将积分化为一个正圆形截面的积分. 这样对于我们取得的闭合曲面$S$上的每个面元$\d S_1$, 总可以取得一个环形管, 找到$S$上的另一个面元与之相对应, 这两个面元的磁通量恰好抵消, 穿过$S$的总磁通恒为$0$, 磁场高斯定理得证.

磁场的高斯定理其实并不要求毕奥-萨伐尔定律的平方反比性. 高斯定理体现的是静磁场的无源性. 安培环路定理要求平方反比性. 安培定理反映的是磁场的有旋性. 两条定理同时使用, 可以完整地反映静磁场的性质.

\section{带电粒子在磁场中的运动}
实验表明, 运动带电粒子在磁场中受到的力有如下关系:
\begin{definition}[洛伦兹力]
  \[
\mathbf{F}=q\mathbf{v}\times \mathbf{B}
\]
\end{definition}

这一力被称为洛伦兹力. 也有定义把运动带电粒子在电场和磁场中受到的合力:

\[
F = q\mathbf{E}+q\mathbf{v}\times\mathbf{B}
\]
称为洛伦兹力. 只要在不引起歧义的情况下, 这两种说法都是可以接受的. 在本文中, 如无特殊说明, 洛伦兹力仅指磁场对带电运动粒子的作用力.

下面我们考察磁场对载流导线的运动. 考虑
图示的这样一段导线.

\insertfig{2-10.png}{洛伦兹力与安培力的关系}{0.25}

微观的角度看, 电流是由导体中的载流子$q$做定向运动而形成的. 设载流子的定向运动速度为$u$, 单位体积内载流子的数目为$n$, 导线的横截面为$S$, 磁感应强度大小为$B$.

先来考察电流大小.$\Delta t$时间内通过某横截面的电量等于对应体积内所有载流子的电量总和:
\[
\Delta Q = n\Delta V =nq Su \Delta t
\]

于是根据定义, 电流强度为通过该截面的电量与截面面积的比值:
\[
I = \frac{\Delta Q}{S}=nqSu
\]
我们考察$\d \mathbf{l}$的一段短导线, 其方向与这段导线内的载流子定向运动速度方向相同. 该段导线所受洛伦兹力总和为:
\[
\d \mathbf{F} = nS\d l q \mathbf{u}\times \mathbf{B} = nqSu \d \mathbf{l}\times \mathbf{B} = I \d \mathbf{l}\times \mathbf{B}
\]

于是我们可以积分得到磁场对于载流导线的作用力为:
\begin{definition}[安培力]
\[
\mathbf{F} =\int_l I \d \mathbf{l} \times \mathbf{B}
\]
\end{definition}

上式被称为安培力公式.

显然, 如果载流回路闭合, 上式积分得$0$. 也就是说对闭合载流回路, 总安培力为$0$.

接下来我们考察磁场对于载流线圈的力矩作用. 我们首先定义线圈的法线方向. 用右手四指环绕方向代表电流的回绕方向, 伸直的拇指即为法线矢量$n$的方向.

\begin{figure}[h]
  \centering
  % 第一张图片
  \begin{minipage}[t]{0.3\textwidth}
    \centering
    \includegraphics[width=\textwidth]{2-11.png}
    \caption{均匀磁场中矩形线圈所受力矩}
  \end{minipage}
  \hfill
  % 第二张图片
  \begin{minipage}[t]{0.3\textwidth}
    \centering
    \includegraphics[width=\textwidth]{2-12.png}
    \caption{左图的投影图}
  \end{minipage}
  \hfill
  \begin{minipage}[t]{0.3\textwidth}
    \centering
    \includegraphics[width=\textwidth]{2-13.png}
    \caption{平面线圈在均匀磁场中受的力矩}
  \end{minipage}
\end{figure}

根据我们的图示, 可以看出$AB$,$CD$两边所受的力对于轴线的力矩为$0$. 对于$BC$,$DA$其受力大小均为$IbB$. 其力臂大小为:
$\frac{a}{2}\sin\theta$. 于是总的力矩大小为:
\[
L = 2IbB\cdot \frac{a}{2}\sin\theta = IabB\sin\theta = ISB\sin\theta = IS (\mathbf{n}\times \mathbf{B})
\]
最后这个结果其实对任意的平面载流线圈都是成立的, 现在我们就来证明这一点.

我们把线圈分割为一系列垂直于转轴的小窄条, 磁场对图示的两个电流元的作用力大小分别为:
\[
\d F_1 = I\d l_1 B\sin\theta_1 \;\d F_1 = I\d l_2 B\sin\theta_2
\]
几何关系告诉我们:
\[
\d l_1 \sin\theta_1 =\d l_2 \sin\theta_2 = \d h
\]
这两个力的合力矩为:
\[
\d L = IB\d h(x_1+x_2) = I B \d S
\]
于是积分可以得到总的力矩仍为$IBS$. 注意到这里我们选取了磁感应强度平行于线圈平面, 当不平行时, 我们只需要使用叉乘将其重新投影到平行方向即可.

在证明过程中我们可以看到, $IS\mathbf{n}$是一个仅与载流线圈本身相关的矢量, 我们定义这个量为线圈的磁矩, 用$m$表示. 于是我们可以把力矩写成:
\[
\mathbf{L} = \mathbf{m} \times \mathbf{B}
\]

\clearpage

